\documentclass[12pt]{article}
\usepackage[english]{babel}
\usepackage{natbib}
\usepackage{url}
\usepackage[utf8]{inputenc}
\usepackage{amsmath}
\usepackage{graphicx}
\graphicspath{{images/}}
\usepackage{parskip}
\usepackage{fancyhdr}
%\usepackage{vmargin}
%\usepackage{dirtytalk}
%\setmarginsrb{3 cm}{2.5 cm}{3 cm}{2.5 cm}{1 cm}{1.5 cm}{1 cm}{1.5 cm}

\title{Grands mécanismes physiopathologiques: les mécanismes endogènes}								% Title
\author{}								% Author
\date{}											% Date

\makeatletter
\let\thetitle\@title
\let\theauthor\@author
\let\thedate\@date
\makeatother

\pagestyle{fancy}
\fancyhf{}
\rhead{\theauthor}
\lhead{\thetitle}
\cfoot{\thepage}

\begin{document}

%%%%%%%%%%%%%%%%%%%%%%%%%%%%%%%%%%%%%%%%%%%%%%%%%%%%%%%%%%%%%%%%%%%%%%%%%%%%%%%%%%%%%%%%%

\begin{titlepage}
	\centering
    \vspace*{0.5 cm}
    \includegraphics[scale = 0.20]{logo.png}\\[1.0 cm]
    \textsc{\LARGE Faculté des Sciences de Montpellier\newline\newline Bioinformatique, Connaissance, Données}\\[2.0 cm]	% University Name
	\textsc{\Large Structures et problématiques de santé}\\[0.5 cm]				% Course Code
	\rule{\linewidth}{0.2 mm} \\[0.4 cm]
	{ \huge \bfseries \thetitle}\\
	\rule{\linewidth}{0.2 mm} \\[1.5 cm]
	

	
	
    
    
    
    
	
\end{titlepage}

%%%%%%%%%%%%%%%%%%%%%%%%%%%%%%%%%%%%%%%%%%%%%%%%%%%%%%%%%%%%%%%%%%%%%%%%%%%%%%%%%%%%%%%%%

\tableofcontents
\pagebreak

%%%%%%%%%%%%%%%%%%%%%%%%%%%%%%%%%%%%%%%%%%%%%%%%%%%%%%%%%%%%%%%%%%%%%%%%%%%%%%%%%%%%%%%%%

\section {Introduction}
études des mécanismes perturbant le fonctionnement physiologique d'un organisme et aboutissant au développement d'une maladie.
Définition de la santé:
\begin{itemize}
\item état de complet  bien être phyique, mental, social, ne consistant pas seulement en une absence de maladie ou d'infirmité.
\item La maladie est une rupture avec la normalité, mais la norme est définie sociologiquement. 
\item Les affections grave dites orgaiques sont au coeur de l'activité médicale traditionnelle.
Pour beaucoup d'autre la limite entre physiologique et pathologiqe n'est pas définie a priori par le biologiste ou le clinicien.
end{itemize}
\end{itemize}
\subsection{Le fonctionnement physiologique n'est pas statique}
Chaque jours au sein d'un organisme, des milliards de cellules:
\begin{itemize}
\item naissent
\item se multiplient, se sdifférencient
\item disparaissent
\end{itemize}

Quotidiennement l'organisme fait fàaceà une multitude de désordres aléatories qu'il neutralise par des mécanismes régulateurs.

\begin{itemize}
\item mutations cellulaires sporadiques, déviances par rapport au développement idéal.
\item foyers de micro-cancérisation
\item réparation de lésions induites par des traumatismes mineurs
\item colonisation invasive ou bactérienne évoluant vers une phase invasive.
\item adaptation aux contraintes de l'environnement
\end{itemize}
Le fonctionnement physiologique suppose: un équilibre dynamique de la croissance cellulaire? une communication régulée entre cellules assurée par le milieu intérieur.

\subsection{Régulation cybernétique de mécanismes physiologiques}
La cybernétique, science des ilotages, peut s'applaiquer à la physiologie. Elle permet une approche non spécialisée mais scientifiquement exacte de la régulation des phénomènes assurant l'homéostasie de l'organisme.

\subsection{La régulation homéostatique un système asservi}

Modalités de la mort cellulaire
\hspace{1 cm}--- Ju

\newpage
\bibliographystyle{plain}
\bibliography{biblist}
enter
\end{document}